\subsection{The targeted halt for every $K$}\label{app:everyK}

Proposition~\ref{prop:target}(ii) is stated for two groups. Its argument needs two populations whose kernels agree to order $(\Delta\eta)^{2K-2}$ and whose largest groups differ at order one. This subsection supplies such a pair for every $K$.

Let $\mu$ be the uniform probability measure on $[-1,1]$, let $p_0,p_1,\ldots$ be its orthonormal polynomials, and put $\lambda_K(s)=1/\sum_{j<K}p_j(s)^2$. Call a rule \emph{admissible} if it has $K$ distinct nodes in $[-1,1]$ and integrates every polynomial of degree at most $2K-2$ exactly against $\mu$. The Gauss rule and the Gauss--Radau rule of Theorem~\ref{thm:kcohort} are admissible.

\begin{lemma}[Largest weights in the family]\label{lem:everyK}
Fix $K\ge2$.
\begin{enumerate}[label=(\roman*),itemsep=1pt]
\item An admissible rule with nodes $t_1,\ldots,t_K$ has weights $\omega_i=\lambda_K(t_i)$. In particular its weights are positive.
\item There is $c_0>0$ such that, for every $|c|\le c_0$, the zeros of $p_K-c\,p_{K-1}$ are the nodes of an admissible rule $R_c$. A point is a node of $R_c$ for at most one $c$.
\item There are $c_1\ne c_2$ in $[-c_0,c_0]$ for which $R_{c_1}$ and $R_{c_2}$ have different largest weights. Their moments of degree $2K-1$ differ.
\end{enumerate}
\end{lemma}

\begin{proof}
(i) Let $\ell_i$ be the Lagrange polynomials of the nodes. Each product $\ell_i\ell_j$ has degree $2K-2$, so the rule integrates it exactly: $\int\ell_i\ell_j\,d\mu=\sum_m\omega_m\ell_i(t_m)\ell_j(t_m)=\omega_i\mathbf 1\{i=j\}$. Taking $i=j$ gives $\omega_i=\int\ell_i^2d\mu>0$. So $\ell_i/\sqrt{\omega_i}$, $i=1,\ldots,K$, is an orthonormal basis of the polynomials of degree at most $K-1$, as is $p_0,\ldots,p_{K-1}$. The sum of squares of an orthonormal basis at a point does not depend on the basis, so $\sum_{j<K}p_j(s)^2=\sum_i\ell_i(s)^2/\omega_i$. At $s=t_i$ the right side is $1/\omega_i$.

(ii) The zeros of $p_K$ are real, simple and inside $(-1,1)$. Simple real zeros move continuously with the coefficients, so for $|c|\le c_0$, with $c_0$ small, $q_c=p_K-c\,p_{K-1}$ has $K$ simple zeros in $(-1,1)$. Give these nodes the interpolatory weights. Let $f$ have degree at most $2K-2$ and write $f=q_cg+\rho$ with $\deg g\le K-2$ and $\deg\rho\le K-1$. Both $p_K$ and $p_{K-1}$ are orthogonal to polynomials of degree at most $K-2$, so $\int q_cg\,d\mu=0$, and the rule integrates $\rho$ exactly and gives $q_cg$ the value zero. Hence the rule is exact for $f$. If $s$ is a zero of $q_c$ and of $q_{c'}$ with $c\ne c'$, then $p_{K-1}(s)=0$ and so $p_K(s)=0$. By the three-term recurrence a common zero of $p_K$ and $p_{K-1}$ is a zero of $p_{K-2}$, and so on down to $p_0=1$, which has none.

(iii) The reciprocal $1/\lambda_K$ is a polynomial of degree $2K-2\ge2$, so every level set of $\lambda_K$ is finite. Suppose all the rules $R_c$, $|c|\le c_0$, had the same largest weight $v$. By (i) each $R_c$ would have a node in the finite set $\{\lambda_K=v\}$, and by (ii) different values of $c$ would give different nodes. There are infinitely many $c$, so this is impossible. For the last statement, let $\nu$ be the difference of the two rules. It is supported on at most $2K$ points. If the two rules also agreed at degree $2K-1$, $\nu$ would vanish on every polynomial of degree at most $2K-1$, including the one that equals $1$ at one support point and $0$ at the others. So $\nu$ would be zero, which contradicts (ii) and the positivity of the weights.
\end{proof}

\begin{proposition}[Supervisory separation for every $K$]\label{prop:targetK}
Fix $K\ge2$ and let $R_{c_1}$ and $R_{c_2}$ be the rules of Lemma~\ref{lem:everyK}(iii), with nodes $s^j_i$, weights $\omega^j_i$ and largest weights $v_1<v_2$. Under the hypotheses of Theorem~\ref{thm:kcohort}, let economy $j\in\{1,2\}$ have $K$ groups with rates $r_c+\Delta\eta\,s^j_i$ and merged weights $W\omega^j_i$, and write $q^j_i=W\omega^j_i/(1-r_c-\Delta\eta\,s^j_i)$.
\begin{enumerate}[label=(\roman*),itemsep=1pt]
\item The response bounds of Theorem~\ref{thm:kcohort} hold for this pair: $\|h^1-h^2\|_2\le C_U(\Delta\eta)^{2K-1}$ for every horizon, and $\|h^1-h^2\|_2\ge\kappa^2W|g|(\Delta\eta)^{2K-1}$ when at least $2K+1$ taps are observed, where $g\ne0$ is the difference of the two rules at degree $2K-1$.
\item The feedback budgets satisfy $|S^1-S^2|=O((\Delta\eta)^{2K-1})$, and the largest groups satisfy $\max_iq^2_i-\max_iq^1_i=W(v_2-v_1)/(1-r_c)+O(\Delta\eta)$.
\item For sufficiently small spacing there is a nonempty interval of ceilings at which halting the largest group meets the ceiling in economy 2 and not in economy 1. For $N$ independent reset records with Gaussian noise $\sigma^2I$, a rule that makes this decision with both errors at most $1/4$ needs
\[
N\ \ge\ \bigl(2\Phi^{-1}(0.75)\,\sigma/\|h^1-h^2\|_2\bigr)^2\ \ge\ c\,\sigma^2(\Delta\eta)^{-(4K-2)} .
\]
\end{enumerate}
\end{proposition}

\begin{proof}
(i) The proof of Theorem~\ref{thm:kcohort} uses four properties of its pair: nodes in $[-1,1]$, positive weights of unit mass, equal moments up to degree $2K-2$, and a nonzero difference at degree $2K-1$. Lemma~\ref{lem:everyK} gives all four. (ii) Expanding $(1-r_c-\Delta\eta\,s)^{-1}$ in powers of $\Delta\eta$, the sum $\sum_iq^j_i$ is $W\sum_{m\ge0}(\Delta\eta)^m\mu^j_m/(1-r_c)^{m+1}$, where $\mu^j_m$ is the moment of degree $m$ of rule $j$. The moments agree for $m\le2K-2$, which gives the first statement. Each $q^j_i$ equals $W\omega^j_i/(1-r_c)+O(\Delta\eta)$, and the maximum is continuous, which gives the second. (iii) The post-halt budgets are $R^j=S^j-\kappa\max_iq^j_i$, so $R^1-R^2=\kappa W(v_2-v_1)/(1-r_c)+O(\Delta\eta)>0$ for small spacing, and any ceiling strictly between $R^2$ and $R^1$ separates the two decisions. The decision is then a test between two Gaussian mean vectors, as in the proof of Proposition~\ref{prop:target}, and (i) gives the bound.
\end{proof}

\paragraph{What is and is not shown.} The pair in Proposition~\ref{prop:targetK} is two members of the family $R_c$, not the Gauss and Gauss--Radau pair itself. For that pair the largest weights differ in every case we computed ($K=2,\ldots,200$), and a proof of that comparison for every $K$ is open. Lemma~\ref{lem:everyK}(ii) proves admissibility only for $c$ near zero. Numerically, for $K=2,\ldots,14$ the rules $R_c$ are admissible on the whole range between the two Gauss--Radau rules, their weights equal $\lambda_K$ at the nodes to $7\times10^{-13}$, and on a grid of 41 values of $c$ the largest weight is smallest at the Gauss rule and largest at the Gauss--Radau rules (\nolinkurl{equilibrium/every_k_halt_check.py}). The counting step in part~(iii) is machine-checked (\texttt{p1\_largest\_weight\_not\_constant} in \nolinkurl{Paper1EveryKHalt.lean}), with parts (i) and (ii) and the finiteness of the level sets as hypotheses. The rest is proved on paper.

